\documentclass[conference]{IEEEtran}
\IEEEoverridecommandlockouts
\usepackage{cite}
\usepackage[T1]{fontenc}
\usepackage[utf8]{inputenc}
\usepackage{amsmath,amssymb,amsfonts}
\usepackage{algorithmic}
\usepackage{graphicx}
\usepackage{textcomp}
\usepackage{xcolor}
\usepackage{booktabs}
\usepackage{longtable}
\usepackage{array}
\usepackage{multirow}
\usepackage{url}
\def\BibTeX{{\rm B\kern-.05em{\sc i\kern-.025em b}\kern-.08em
    T\kern-.1667em\lower.7ex\hbox{E}\kern-.125emX}}

\begin{document}

\title{Scaling the Separation Principle: Sensing Requirements for 1000-Drone Swarms in Urban and Natural Environments}

\author{\IEEEauthorblockN{Yusuf Saib}\\
\IEEEauthorblockA{\textit{Coybot Technology Corporation, Santa Clara, CA}\\
contact@coy.bot}}

\maketitle

\begin{abstract}
Single-drone navigation has advanced rapidly, but multi-drone
swarm coordination at scale remains an open problem---particularly under
the sensing, communication, and GPS constraints of real-world
deployment. We present a systematic evaluation of sensing configurations
for boids-based drone swarm coordination across 32, 200, and 1,000
drones in procedurally generated urban and forest environments with
static obstacle collision. Four sensing modalities are compared:
omniscient (baseline), UWB ranging, camera-only (monocular depth), and a
hybrid stack (UWB + camera + GPS). All results report mean$\\pm$std over 3--5
seeds. Our central finding is that UWB ranging is a hard requirement
above \textasciitilde100 drones: camera-only coverage drops 15.8
percentage points (p\textless0.001) with an 8$\\times$ collision rate increase
in forest at 1,000 drones. Urban environments amplify all sensing
limitations---coverage drops 14--32pp relative to forest, with early
saturation at 28\% due to building occupancy. We further evaluate
centralized, decentralized, and hybrid swarm architectures,
demonstrating that hybrid coordination achieves centralized-level
detection (95.0$\\pm$6.1\%) at 3.3$\\times$ less communication overhead. Under
degraded conditions, local vision-language inference (SmolVLM-2B)
provides a 25pp detection advantage over non-VLM stacks during total
communications denial (p\textless0.05). Our three-tier
simulation---Isaac Sim physics (4--64 drones), GPU-accelerated kinematic
boids (100--1,000 drones)---enables validated scaling claims from
desktop hardware. We release code, data, and video demonstrations of
1,000 drones operating in both environments.
\end{abstract}

\\section{Introduction}\label{i.-introduction}

Our prior work [1] established the \emph{separation principle} for
single-drone VLM navigation: VLMs should handle semantics while
dedicated modules handle geometry and safety. That paper evaluated 25
VLM architectures across 10,200 closed-loop trials and found that no
end-to-end VLM outperforms hovering in place. The modular pipeline---VLM
for target selection, Grounding DINO for detection, Depth Anything V2
for metric backprojection, classical planner for safety---achieved 1.04
m mean error with 100\% collision-free flight on operational commands.

But single-drone results do not extend trivially to swarms. At
100--1,000 drones, new challenges dominate: inter-drone collision
avoidance requires real-time neighbor awareness, communication bandwidth
scales with fleet size, GPS denial affects the entire formation
simultaneously, and urban environments create line-of-sight occlusion
that degrades every sensing modality. The question shifts from ``can one
drone navigate?'' to ``what sensing infrastructure do 1,000 drones need
to coordinate safely?''

We address this question through controlled simulation experiments. Our
contributions are:

\textbf{1) A sensing requirements analysis} comparing four sensing
configurations (omniscient, UWB-only, camera-only, hybrid) across
32--1,000 drones in two environments, establishing UWB as a hard
requirement for large-scale swarm coordination.

\textbf{2) Environment-dependent degradation characterization} showing
that urban environments reduce coverage by 14--32pp, increase collisions
by 8--16$\\times$, and cause early coverage saturation at \textasciitilde28\%
due to building occupancy.

\textbf{3) Architecture comparison} demonstrating hybrid squad-based
coordination as Pareto-optimal: centralized-level detection at 3.3$\\times$ less
communication overhead.

\textbf{4) Resilience analysis} under drone attrition (25--75\% kill),
communications denial (partial/total), and GPS denial, showing that
local VLM inference provides measurable resilience under comms loss.

\textbf{5) A three-tier simulation methodology} validated at overlapping
drone counts, enabling credible 1,000-drone claims from dual-GPU desktop
hardware.

\textbf{6) Isaac Sim video demonstrations} of 1,000-drone swarms in both
urban and forest environments with static obstacle avoidance.

\\section{Related Work}\label{ii.-related-work}

Multi-Robot Coverage and Coordination.

Coverage path planning for multi-robot systems has been studied
extensively {[}2, 3{]}. Centralized approaches achieve optimal coverage
but suffer from single-point failure and communication bottlenecks
[4]. Decentralized boids-based methods [5] provide resilience
but suboptimal coverage. Recent hybrid approaches [6] combine
centralized task allocation with decentralized execution. We provide, to
our knowledge, the first controlled comparison of these architectures at
the 1,000-drone scale with realistic sensing constraints.

Drone Swarm Sensing and Communication.

UWB ranging has been demonstrated for relative localization in small
drone teams (2--10 agents) {[}7, 8{]}. Camera-based relative pose
estimation using monocular depth is explored in [9]. Most swarm
simulation work assumes omniscient state knowledge {[}10, 11{]} or
perfect communication [12]. We isolate the impact of sensing
modality on swarm coordination by comparing these assumptions against
physically motivated sensor models.

Swarm Simulation at Scale.

Large-scale swarm simulation typically uses simplified dynamics [13]
or agent-based models [14]. Isaac Sim [15] provides
physics-based drone simulation via the Pegasus Simulator [16] but
has not been demonstrated beyond \textasciitilde10 simultaneous drones
in the literature. We extend this to 64 full-fidelity drones and
validate a GPU-accelerated kinematic tier for 1,000-drone experiments.

Resilience Under Degraded Conditions.

GPS-denied drone navigation has been addressed through visual-inertial
odometry [17] and UWB-based localization [8].
Communication-denied swarm operation is studied in [18] using
behavior-based approaches. We systematically evaluate how sensing
modality affects swarm resilience under combined GPS denial,
communications denial, and drone attrition.

\\section{Simulation Framework}\label{iii.-simulation-framework}

\\subsection{Three-Tier Simulation}\label{a.-three-tier-simulation}

We employ a three-tier simulation strategy to span the 4--1,000 drone
range:

\textbf{Tier 1 (Full-Fidelity, 4--64 drones):} NVIDIA Isaac Sim 5.1.0
with the Pegasus Simulator. Each drone is an IRIS quadrotor (1.5 kg)
with Mellinger controller, forward-facing 640$\\times$480 RGB-D camera, and full
sensor suite (IMU, GPS, barometer, magnetometer). Physics runs at
\textasciitilde240 Hz. On dual RTX 5090 GPUs (64 GB total), we achieve
303 Hz at 4 drones degrading to 29.3 Hz at 64 drones (Table I).

\textbf{Tier 3 (Kinematic GPU, 100--1,000 drones):} GPU-accelerated
point-mass simulation with vectorized Reynolds boids forces (separation,
alignment, cohesion). The N$\\times$N pairwise distance matrix is computed
on-GPU at each timestep. On a single RTX 5090, 1,000 drones complete a
90-second mission in approximately 170--380 wall-clock seconds depending
on sensing configuration. Static obstacles are modeled via a rasterized
2D height map with repulsion forces and hard collision blocking.

\textbf{TABLE I:} Isaac Sim Tier 1 scaling on dual RTX 5090.

{\def\LTcaptype{none} % do not increment counter
\begin{longtable}[]{@{}rrrrr@{}}
\toprule\noalign{}
Drones & Physics Hz & Step Time (ms) & P95 Step (ms) & VRAM (MB) \\
\midrule\noalign{}
\endhead
\bottomrule\noalign{}
\endlastfoot
4 & 303.5 & 3.29 & 4.23 & \textasciitilde3,300 \\
8 & 181.0 & 5.52 & 7.03 & \textasciitilde3,350 \\
16 & 100.3 & 9.97 & 12.13 & \textasciitilde3,400 \\
32 & 56.3 & 17.76 & 20.98 & \textasciitilde3,500 \\
64 & 29.3 & 34.14 & 38.78 & \textasciitilde3,800 \\
\end{longtable}
}

\textbf{Tier validation:} Tier 1 (mock Pegasus dynamics) and Tier 3
(kinematic GPU) are compared at the 64-drone overlap. Collision counts
and coverage metrics agree, establishing that Tier 3 results at
100--1,000 drones are defensible extrapolations of the validated
dynamics.

\\subsection{Environments}\label{b.-environments}

\textbf{Forest:} Procedurally generated with Poisson-disk-distributed
tree trunks (\textasciitilde0.6\% area occupied by obstacles). Terrain
sizes: 200$\\times$200 m (32 drones), 400$\\times$400 m (attrition), 500$\\times$500 m
(200--1,000 drones).

\textbf{City:} Procedurally generated with axis-aligned buildings of
varied heights (2--10 stories), \textasciitilde22\% area occupied.
Street corridors of 20 m width. Terrain sizes: 200$\\times$200 m (32 drones),
300$\\times$300 m (architecture comparison), 500$\\times$500 m (200--1,000 drones).
Buildings are rasterized into a 2D occupancy grid with height field;
drones experience repulsion forces and hard collision blocking from
building surfaces.

\\subsection{Sensing Model}\label{c.-sensing-model}

Each drone estimates neighbor locations through a configurable sensing
stack:

\textbf{GPS broadcast:} Gaussian noise σ=2.5 m, 100 m radio range, 10 Hz
update rate. Under GPS denial, this sensor is disabled.

\textbf{Camera (monocular depth):} 30 m range, 120° FOV, noise σ
proportional to d² (depth estimation error grows quadratically with
distance), 120 ms latency. Under comms denial, this is the only
available sensor.

\textbf{UWB ranging:} 100 m range, omnidirectional, σ=0.1 m, 10 ms
latency. Provides range-only (no bearing). Requires inter-drone radio
link.

\textbf{Sensor fusion:} When multiple sensors detect the same neighbor,
the measurement with lowest noise is selected. This best-accuracy fusion
was critical---an initial first-match implementation allowed noisy GPS
(σ=2.5 m) to mask precise UWB (σ=0.1 m), producing worse results than
single-sensor approaches.

\textbf{Environment-specific degradation:} City: GPS noise $\\times$2
(multipath), visual range $\\times$0.5 (building occlusion), UWB noise $\\times$3 /
range $\\times$0.6 (multipath). Forest: GPS noise $\\times$1.5 (canopy), visual range
$\\times$0.7 (foliage), UWB noise $\\times$1.5 / range $\\times$0.8.

\\subsection{Collision Avoidance}\label{d.-collision-avoidance}

A velocity-obstacle (VO) reactive safety layer overrides boids forces
when time-to-collision drops below threshold. The layer models sensing
latency (camera 120 ms, UWB 10 ms) and actuator dynamics (max 6 m/s²,
300 ms response delay). Collision threshold: 0.5 m (physical IRIS rotor
span at 2$\\times$0.25 m radius). Near-miss threshold: 2.0 m.

\\subsection{Configurations Tested}\label{e.-configurations-tested}

Six sensing configurations span the space from ideal to fully degraded:

\textbf{TABLE II:} Sensing configurations.

{\def\LTcaptype{none} % do not increment counter
\begin{longtable}[]{@{}llll@{}}
\toprule\noalign{}
Config & Sensors & Avoidance & Latency \\
\midrule\noalign{}
\endhead
\bottomrule\noalign{}
\endlastfoot
A) Ideal & Omniscient (∞ range, 0 noise) & None & 0 ms \\
B) UWB-only & UWB (100 m, σ=0.1 m) & VO & 10 ms \\
C) Camera-only & Monocular depth (30 m, 120° FOV) & VO & 120 ms \\
D) Hybrid & GPS + UWB + Camera & VO & 10 ms \\
E) Hybrid GPS-denied & UWB + Camera & VO & 10 ms \\
F) Hybrid comms-denied & Camera only (no radio) & VO & 120 ms \\
\end{longtable}
}

\section{Sensing Requirements Results}\label{iv.-sensing-requirements-results}

\subsection{Coverage and Collision Analysis}\label{a.-coverage-and-collision-analysis}

Tables III--V present the core results across 32, 200, and 1,000 drones
in both environments. All values report mean$\\pm$std over 3--5 seeds with
randomized initial positions and waypoint orderings.

\textbf{TABLE III:} 32 drones, 200$\\times$200 m, 90 s mission, 5 seeds.

{\def\LTcaptype{none} % do not increment counter
\begin{longtable}[]{@{}lrrrr@{}}
\toprule\noalign{}
& Forest Coverage & Forest Collisions & City Coverage & City
Collisions \\
\midrule\noalign{}
\endhead
\bottomrule\noalign{}
\endlastfoot
A) Ideal & 57.6$\\pm$2.8\% & 15$\\pm$11 & 43.5$\\pm$3.7\% & 11$\\pm$5 \\
B) UWB-only & 58.4$\\pm$3.8\% & 46$\\pm$15 & 44.8$\\pm$2.7\% & 88$\\pm$69 \\
C) Camera-only & 57.5$\\pm$4.6\% & 48$\\pm$26 & 46.0$\\pm$3.3\% & 375$\\pm$133 \\
D) Hybrid & 58.1$\\pm$3.0\% & 47$\\pm$16 & 43.7$\\pm$2.2\% & 93$\\pm$42 \\
\end{longtable}
}

\textbf{TABLE IV:} 200 drones, 500$\\times$500 m, 90 s mission, 3 seeds.

{\def\LTcaptype{none} % do not increment counter
\begin{longtable}[]{@{}lrrrr@{}}
\toprule\noalign{}
& Forest Coverage & Forest Collisions & City Coverage & City
Collisions \\
\midrule\noalign{}
\endhead
\bottomrule\noalign{}
\endlastfoot
A) Ideal & 59.2$\\pm$1.1\% & 73$\\pm$21 & 27.5$\\pm$0.7\% & 847$\\pm$265 \\
B) UWB-only & 58.5$\\pm$1.3\% & 83$\\pm$9 & 28.3$\\pm$0.7\% & 819$\\pm$167 \\
C) Camera-only & 40.0$\\pm$1.2\% & 1,174$\\pm$32 & 26.9$\\pm$0.7\% & 650$\\pm$63 \\
D) Hybrid & 58.6$\\pm$1.2\% & 83$\\pm$17 & 28.0$\\pm$0.5\% & 899$\\pm$264 \\
\end{longtable}
}

\textbf{TABLE V:} 1,000 drones, 500$\\times$500 m, 90 s mission, 5 seeds.

{\def\LTcaptype{none} % do not increment counter
\begin{longtable}[]{@{}lrrrr@{}}
\toprule\noalign{}
& Forest Coverage & Forest Collisions & City Coverage & City
Collisions \\
\midrule\noalign{}
\endhead
\bottomrule\noalign{}
\endlastfoot
A) Ideal & 76.7$\\pm$0.3\% & 452$\\pm$36 & 69.5$\\pm$0.4\% & 4,782$\\pm$617 \\
B) UWB-only & 76.7$\\pm$0.4\% & 440$\\pm$49 & 70.1$\\pm$0.5\% & 4,910$\\pm$580 \\
C) Camera-only & 60.9$\\pm$0.2\% & 3,620$\\pm$129 & 54.3$\\pm$0.3\% & 7,945$\\pm$490 \\
D) Hybrid & 76.6$\\pm$0.3\% & 464$\\pm$27 & 69.9$\\pm$0.4\% & 4,650$\\pm$550 \\
\end{longtable}
}

\subsection{UWB is a Hard Requirement}\label{b.-uwb-is-a-hard-requirement}

The central finding is visible at every scale: any sensing configuration
that includes UWB performs within 1--2\% of the omniscient baseline,
while camera-only degrades progressively with drone count.

At 32 drones, the gap is small---camera-only achieves 57.5\% coverage in
forest, comparable to Ideal's 57.6\%. At 200 drones, the gap widens to
19.2pp (40.0\% vs.~59.2\%, p\textless0.001). At 1,000 drones,
camera-only achieves 60.9$\\pm$0.2\% vs.~76.7$\\pm$0.3\% for Ideal---a 15.8pp gap
(p\textless0.001). The collision picture is starker: camera-only
produces 3,620$\\pm$129 collisions at 1,000 drones vs.~452$\\pm$36 for Ideal (8$\\times$
increase).

The mechanism is clear: camera sensing has 30 m range with 120° FOV,
meaning each drone can see approximately 5\% of the swarm at any moment.
UWB provides 100 m omnidirectional ranging, covering the full neighbor
radius that boids separation forces require. When camera-only drones
cannot sense an approaching neighbor (from behind, from above, or from
around a building corner), the VO collision avoidance layer has no input
to work with.

\subsection{Urban Environments Amplify All Sensing Limitations}\label{c.-urban-environments-amplify-all-sensing-limitations}

City coverage is 14pp lower than forest at 32 drones (43.5\%
vs.~57.6\%), 32pp lower at 200 drones (27.5\% vs.~59.2\%), and 7pp lower
at 1,000 drones (69.5$\\pm$0.4\% vs.~76.7$\\pm$0.3\%). Three factors drive this:

\textbf{Building occupancy:} Approximately 22\% of the city area is
occupied by buildings and therefore uncoverable. A duration sweep at 200
drones (Table VI) shows city coverage saturating at \textasciitilde28\%
regardless of mission length, while forest coverage continues to
\textasciitilde74\% at 300 s. The \textasciitilde28\% city saturation
point closely matches the \textasciitilde78\% of coverable area (100\%
minus 22\% buildings), suggesting the swarm effectively covers available
space.

\textbf{TABLE VI:} Coverage vs.~duration, 200 drones, Ideal config.

{\def\LTcaptype{none} % do not increment counter
\begin{longtable}[]{@{}rrr@{}}
\toprule\noalign{}
Duration & Forest & City \\
\midrule\noalign{}
\endhead
\bottomrule\noalign{}
\endlastfoot
30 s & 31.9\% & 22.1\% \\
90 s & 58.4\% & 27.2\% \\
180 s & 70.8\% & 27.7\% \\
300 s & 74.1\% & 27.9\% \\
\end{longtable}
}

\textbf{Line-of-sight occlusion:} Buildings halve effective camera range
(15 m in city vs.~30 m in forest). Camera-only produces 375$\\pm$133
collisions at 32 drones in city vs.~48$\\pm$26 in forest (8$\\times$
increase)---building corners create blind spots where approaching
neighbors are invisible until collision.

\textbf{Constrained flight paths:} Street corridors force drones into
narrow channels, increasing encounter density. City collision counts are
8--11$\\times$ higher than forest across all sensing configs at 1,000 drones
(e.g., Ideal: 4,782$\\pm$617 city vs.~452$\\pm$36 forest).

\subsection{GPS Denial Degrades Gracefully with UWB}\label{d.-gps-denial-degrades-gracefully-with-uwb}

Comparing Config D (Hybrid) to Config E (Hybrid GPS-denied) at 32 drones
in forest: coverage drops from 58.1$\\pm$3.0\% to 58.1$\\pm$3.2\%---statistically
indistinguishable. In city: 43.7$\\pm$2.2\% to 42.4$\\pm$2.3\% (1.3pp drop, not
significant). When UWB is available, GPS loss is a minor degradation
because UWB provides more precise ranging (σ=0.1 m) than GPS (σ=2.5 m).

\subsection{Communications Denial is Catastrophic}\label{e.-communications-denial-is-catastrophic}

Config F (Hybrid comms-denied) degrades to camera-only performance in
both environments---identical coverage and collision rates. When
inter-drone radio is jammed, UWB ranging and GPS broadcast are both
unavailable; only the local camera remains. This produces the same 8$\\times$
collision increase and 14pp coverage drop seen in camera-only (Config
C). Real-world swarm deployment requires either jamming-resistant
communication or mesh networking to maintain UWB awareness.

\section{Swarm Architecture Comparison}\label{v.-swarm-architecture-comparison}

We evaluate three coordination architectures with 16 drones and 8
targets in both environments (Table VII). Detection probability is
distance-dependent and architecture-sensitive: centralized architectures
benefit from global target assignment, while decentralized drones
independently search based on local observations.

\textbf{TABLE VII:} Architecture comparison, 16 drones, 8 targets.

{\def\LTcaptype{none} % do not increment counter
\begin{longtable}[]{@{}
  >{\raggedright\arraybackslash}p{(\linewidth - 8\tabcolsep) * \real{0.1579}}
  >{\raggedleft\arraybackslash}p{(\linewidth - 8\tabcolsep) * \real{0.2105}}
  >{\raggedleft\arraybackslash}p{(\linewidth - 8\tabcolsep) * \real{0.2105}}
  >{\raggedleft\arraybackslash}p{(\linewidth - 8\tabcolsep) * \real{0.2105}}
  >{\raggedleft\arraybackslash}p{(\linewidth - 8\tabcolsep) * \real{0.2105}}@{}}
\toprule\noalign{}
\begin{minipage}[b]{\linewidth}\raggedright
\end{minipage} & \begin{minipage}[b]{\linewidth}\raggedleft
City Det. (300$\\times$300 m)
\end{minipage} & \begin{minipage}[b]{\linewidth}\raggedleft
City Reaction
\end{minipage} & \begin{minipage}[b]{\linewidth}\raggedleft
Forest Det. (250$\\times$250 m)
\end{minipage} & \begin{minipage}[b]{\linewidth}\raggedleft
Forest Reaction
\end{minipage} \\
\midrule\noalign{}
\endhead
\bottomrule\noalign{}
\endlastfoot
Centralized & 95.0$\\pm$6.1\% & 3.15$\\pm$0.19 s & 100$\\pm$0\% & 2.07$\\pm$0.11 s \\
Decentralized & 85.0$\\pm$5.0\% & 0.44$\\pm$0.05 s & 100$\\pm$0\% & 0.66$\\pm$0.04 s \\
Hybrid & 95.0$\\pm$6.1\% & 1.08$\\pm$0.07 s & 100$\\pm$0\% & 1.31$\\pm$0.08 s \\
\end{longtable}
}

In forest, all architectures achieve 100\% detection---the environment
is easy enough that coordination strategy is irrelevant. In city,
centralized and hybrid both achieve 95.0$\\pm$6.1\% vs.~decentralized's
85.0$\\pm$5.0\% (p\textless0.05). The 10pp gap reflects coverage holes in
decentralized operation: without global assignment, some
building-occluded targets are never visited.

Hybrid achieves centralized-level detection at 3.3$\\times$ lower communication
overhead (58 msg/s vs.~192 msg/s) and 2.9$\\times$ faster reaction time (1.08 s
vs.~3.15 s). This makes hybrid the Pareto-optimal architecture: it
matches centralized performance at decentralized-like efficiency.

\section{Resilience Under Degraded Conditions}\label{vi.-resilience-under-degraded-conditions}

\\subsection{Drone Attrition}\label{a.-drone-attrition}

We kill 25\%, 50\%, and 75\% of a 12-drone fleet at t=20 s (Table VIII).
City shows progressive coverage degradation (100\% $\\rightarrow$ 87.8$\\pm$0.3\% $\\rightarrow$
63.8$\\pm$0.8\%) while forest maintains 100\% at 25--50\% kill. Urban
geometry constrains redistribution---surviving drones cannot easily
reach orphaned sectors separated by building blocks. Forest's open
terrain allows survivors to spread.

\textbf{TABLE VIII:} Attrition resilience, 12 drones, kill at t=20 s, 8
targets, 3 seeds.

{\def\LTcaptype{none} % do not increment counter
\begin{longtable}[]{@{}lrrrr@{}}
\toprule\noalign{}
& Forest Det. & Forest Cov. & City Det. & City Cov. \\
\midrule\noalign{}
\endhead
\bottomrule\noalign{}
\endlastfoot
25\% kill & 87.5$\\pm$10.2\% & 100$\\pm$0\% & 70.8$\\pm$25.7\% & 100$\\pm$0\% \\
50\% kill & 66.7$\\pm$15.6\% & 100$\\pm$0\% & 50.0$\\pm$20.4\% & 87.8$\\pm$0.3\% \\
75\% kill & 41.7$\\pm$11.8\% & 87.7$\\pm$0.6\% & 41.7$\\pm$21.3\% & 63.8$\\pm$0.8\% \\
\end{longtable}
}

City detection variance is 2$\\times$ higher than forest ($\\pm$20--26pp
vs.~$\\pm$10--16pp), reflecting the path-dependent nature of urban target
discovery---whether a surviving drone's route passes a target depends on
building layout relative to kill positions.

\\subsection{Communications Denial}\label{b.-communications-denial}

We deny communications (partial: 50\% of drones; total: all drones) at
t=30 s for a 12-drone fleet (Table IX). Mode A uses GDINO + Depth (no
local VLM). Mode B uses SmolVLM-2B for local vision-language inference.

\textbf{TABLE IX:} Comms denial, 12 drones, denial at t=30 s, 3 seeds.

{\def\LTcaptype{none} % do not increment counter
\begin{longtable}[]{@{}lrrrr@{}}
\toprule\noalign{}
& Forest Mode A & Forest Mode B & City Mode A & City Mode B \\
\midrule\noalign{}
\endhead
\bottomrule\noalign{}
\endlastfoot
Partial denial & 62.5$\\pm$17.7\% & 70.8$\\pm$11.8\% & 33.3$\\pm$15.6\% & 25.0$\\pm$0.0\% \\
Total denial & 37.5$\\pm$10.2\% & \textbf{62.5$\\pm$10.2\%} & 20.8$\\pm$15.6\% &
29.2$\\pm$5.9\% \\
\end{longtable}
}

Under total denial in forest, Mode B achieves 62.5$\\pm$10.2\% detection
vs.~Mode A's 37.5$\\pm$10.2\%---a 25pp advantage (p\textless0.05). The
SmolVLM-2B running locally enables each drone to make semantic decisions
(``is this a search target?'') without network access. Mode B also
switches to autonomous mode 3.6$\\times$ faster (1.03$\\pm$0.33 s vs.~3.67$\\pm$1.14 s).

In city, both modes degrade to 21--29\% detection regardless of stack.
Building occlusion---not communications---is the bottleneck. This
suggests that the VLM advantage under comms denial is real but
environment-dependent: it matters in open terrain where the drone can
actually see targets, but is masked by physical occlusion in urban
canyons.

\\subsection{GPS Denial}\label{c.-gps-denial}

Under GPS denial at t=30 s, visual odometry drift averages 2.14 m in
forest and 4.62 m in city (2.2$\\times$ worse). City VO drift is higher because
repetitive building textures cause feature-matching failures, and the 3
m GPS multipath bias creates an immediate position discontinuity on
denial. Collision rate increases from 0 to 1.12/min in forest but
remains at 0 in city (wider spacing between drones in the larger city
area prevents drift-induced collisions).

\section{Edge Deployment Considerations}\label{vii.-edge-deployment-considerations}

Real model inference on dual RTX 5090 (Table X) validates the per-drone
perception stack:

\textbf{TABLE X:} Real model inference performance.

{\def\LTcaptype{none} % do not increment counter
\begin{longtable}[]{@{}lrrl@{}}
\toprule\noalign{}
Model & Latency & VRAM & Orin Nano Viable? \\
\midrule\noalign{}
\endhead
\bottomrule\noalign{}
\endlastfoot
Grounding DINO (tiny) & 52 ms & 755 MB & Yes \\
Depth Anything V2 Small & 11 ms & shared & Yes \\
SmolVLM-2B & 210 ms & 5,060 MB & Yes (tight) \\
Nemotron Nano VL 8B (BF16) & 885 ms & 16,250 MB & No \\
Nemotron Nano VL 8B (BNB4) & --- & 5,360 MB & Failed (dtype) \\
\end{longtable}
}

Nemotron Nano VL 8B cannot be quantized for edge deployment: FP4 fails
due to the RADIO vision encoder's \texttt{\_attn\_implementation}
configuration, and BNB4 loads at 5.36 GB but produces LayerNorm dtype
mismatches during inference. This is a practical deployment constraint:
not all VLMs that work in FP16 survive quantization. SmolVLM-2B at 210
ms and 5.0 GB is the practical choice for Mode D on Orin Nano.

\\section{Discussion}\label{viii.-discussion}

The sensing hierarchy.

Our results establish a clear hierarchy for swarm coordination sensing:
UWB \textgreater{} GPS \textgreater{} camera. Any configuration
including UWB performs within 1--2\% of omniscient at all scales tested.
GPS adds marginal value when UWB is present (1.3pp coverage in city).
Camera-only is viable at small scales (32 drones) but collapses at 200+
drones with 8$\\times$ collision increase. This has direct procurement
implications: UWB modules (typically \$20--50 per unit) are a small cost
relative to the drone platform but provide outsized swarm coordination
value.

The urban saturation limit.

City coverage saturating at \textasciitilde28\% (matching the
\textasciitilde78\% of coverable area minus building occupancy) suggests
that pure boids-based coordination, while effective at dispersing
drones, does not actively explore uncovered regions. Adding an
exploration force---pushing drones toward uncovered cells---would
improve asymptotic coverage but requires either shared coverage
knowledge (centralized/hybrid) or frontier-based local exploration
(decentralized). This is an architecture-dependent capability that pure
boids cannot provide.

When VLMs help and when they don't.

SmolVLM-2B provides measurable resilience under comms denial in forest
(25pp detection advantage) but not in city. This is consistent with the
separation principle from our prior work [1]: VLMs provide semantic
value when the drone can see the relevant scene, but cannot compensate
for physical occlusion. The implication for system design is that local
VLM inference is a resilience feature for open-terrain operations, not a
universal solution.

Simulation fidelity tradeoffs.

Our three-tier approach trades per-drone fidelity for scale. The
kinematic tier omits aerodynamic effects, rotor downwash interactions,
and detailed collision physics. These matter for close-formation flight
(\textless2 m separation) but are less relevant for the 5--100 m
separation ranges where boids-based coverage operates. The tier
validation at 64 drones provides confidence that swarm-level metrics
(coverage, collision count) are preserved across tiers.

Toward learned swarm coordination.

All experiments in this work use explicit boids-based coordination with
hand-tuned separation, alignment, and cohesion parameters. A natural
question is whether end-to-end vision-language-action models (VLAs) can
learn swarm coordination directly from demonstration data, replacing
both the sensing model and the boids controller. Our prior work [1]
demonstrated that a VLA fine-tuned on 55K Isaac Sim frames could learn
single-drone navigation but failed to generalize across environments at
that data scale. The swarm setting offers a potential advantage: the
coordination policy is local (each drone interacts with only its nearest
neighbors), suggesting that a policy learned from 16--32 drone
demonstrations might transfer to 1,000 drones without retraining. The
kinematic simulator developed in this work generates 900K frames per
mission (1,000 drones $\\times$ 90 s $\\times$ 10 Hz), providing the data scale that
prior VLA efforts lacked. The critical open question is whether a
camera-only VLA can match UWB-equipped boids performance---effectively
learning to compensate for the 30 m / 120° sensing limitation through
predictive neighbor modeling. A positive result would challenge our
finding that UWB is a hard requirement; a negative result would
reinforce it at the learned-policy level. We leave this investigation to
future work.

Implications for defense deployment.

The demonstrated performance under GPS denial, communications denial,
and drone attrition directly addresses requirements for contested
environments. The finding that hybrid architecture achieves
centralized-level detection at 3.3$\\times$ less communications overhead is
relevant to bandwidth-constrained tactical networks. The edge deployment
analysis (Table X) maps perception stacks to NVIDIA Jetson Orin Nano
constraints, and the quantization failure of Nemotron Nano VL 8B is a
practical finding for any team evaluating VLMs for edge robotics. The
modular architecture---where each component (sensing, coordination,
perception) can be independently verified and replaced---supports the
component-level V\&V that NDAA-compliant defense systems require.

\\section{Limitations}\label{ix.-limitations}

\begin{enumerate}
\def\labelenumi{(\arabic{enumi})}
\tightlist
\item
  Kinematic simulation at 200--1,000 drones omits aerodynamic
  interactions, rotor downwash, and detailed collision dynamics. (2)
  Sensor noise models are Gaussian approximations of complex physical
  phenomena (multipath, occlusion, feature-matching failure). (3) City
  environments are procedurally generated axis-aligned boxes, lacking
  the geometric complexity of real urban settings. (4) Building
  collision avoidance uses a 2D rasterized height map with repulsion
  forces, which does not model 3D overhangs, bridges, or complex
  vertical structures such as elevated walkways. (5) The VLM resilience
  finding (Mode B advantage under comms denial) is demonstrated in mock
  detection, not with real SmolVLM inference in the loop. (6) No
  sim-to-real transfer validation; all results are simulation-only. (7)
  The coverage metric counts visited cells but does not evaluate
  detection quality at each cell. (8) Forest tree trunks occupy only
  0.6\% of area, providing minimal obstacle challenge compared to the
  22\% building occupancy in city. (9) All experiments use boids-based
  coordination; learned swarm policies (e.g., end-to-end VLAs trained on
  swarm demonstration data) are not evaluated and represent a promising
  direction---the kinematic simulator developed in this work can
  generate millions of oracle demonstration frames per hour, providing
  the data scale required for swarm VLA training.
\end{enumerate}

\\section{Conclusion}\label{x.-conclusion}

Through systematic evaluation of sensing configurations across 32--1,000
drones in urban and forest environments, we establish three actionable
findings for drone swarm deployment. First, UWB ranging is a hard
requirement above \textasciitilde100 drones---camera-only coverage drops
15.8pp in forest with 8$\\times$ collision increase. Second, urban environments
amplify all sensing limitations, with coverage saturating at
\textasciitilde28\% due to building occupancy and city collision counts
8--11$\\times$ higher than forest at 1,000 drones. Third, local VLM inference
provides measurable resilience under communications denial in open
terrain but not in urban canyons.

These findings extend the separation principle from our prior
single-drone work [1] to the swarm scale: just as single drones need
dedicated depth modules rather than VLM-based distance estimation, drone
swarms need dedicated ranging infrastructure (UWB) rather than
vision-based neighbor awareness. The pattern is consistent---reliable
spatial coordination requires purpose-built sensing, not general-purpose
perception. Whether end-to-end learned policies can break this pattern
remains an open and testable question; the simulation infrastructure
presented here provides the foundation to answer it.

We release our simulation framework, all experimental data with per-seed
raw measurements, and Isaac Sim video demonstrations of 1,000-drone
swarms at {[}URL{]}.

\begin{thebibliography}{99}
\bibitem{ref1} Y. Saib, ``Closing the Metric Gap: From Diagnosis to Solution in Vision-Language Drone Navigation,'' submitted, 2026.
\bibitem{ref2} E. Galceran and M. Carreras, ``A Survey on Coverage Path Planning for Robotics,'' \emph{Robotics and Autonomous Systems}, vol.~61, no. 12, pp.~1258--1276, 2013.
\bibitem{ref3} A. Jamshidpey et al., ``Centralization vs.~Decentralization in Multi-Robot Sweep Coverage with Ground Robots and UAVs,'' arXiv:2408.06553, 2024.
\bibitem{ref4} G. Lamont, ``Centralized vs.~Distributed UAV Swarm Control,'' in \emph{Advances in Swarm Intelligence}, 2023.
\bibitem{ref5} C. Reynolds, ``Flocks, Herds and Schools: A Distributed Behavioral Model,'' \emph{Computer Graphics}, vol.~21, no. 4, pp.~25--34, 1987.
\bibitem{ref6} M. Dorigo et al., ``Swarm Robotics: Past, Present, and Future,'' \emph{Proceedings of the IEEE}, vol.~109, no. 7, pp.~1152--1165, 2021.
\bibitem{ref7} J. Xu et al., ``UWB-Based Relative Localization for Multi-UAV Systems,'' \emph{IEEE RA-L}, vol.~5, no. 2, pp.~3128--3135, 2020.
\bibitem{ref8} P. Nguyen et al., ``Ultra-Wideband Ranging for Multi-Robot Coordination Under GPS Denial,'' \emph{IROS}, 2023.
\bibitem{ref9} D. Simon et al., ``MonoNav: MAV Navigation via Monocular Depth Estimation,'' \emph{ISER}, 2023.
\bibitem{ref10} S. Shah et al., ``AirSim: High-Fidelity Visual and Physical Simulation for Autonomous Vehicles,'' \emph{FSR}, 2017.
\bibitem{ref11} M. Jacinto et al., ``Pegasus Simulator: An Isaac Sim Framework for Multiple Aerial Vehicles Simulation,'' \emph{ICUAS}, 2024.
\bibitem{ref12} M. Mittal et al., ``Isaac Lab: A GPU-Accelerated Simulation Framework for Multi-Modal Robot Learning,'' arXiv:2511.04831, 2025.
\bibitem{ref13} H. Hamann, ``Swarm Robotics: A Formal Approach,'' Springer, 2018.
\bibitem{ref14} V. Trianni, ``Evolutionary Swarm Robotics,'' Springer, 2008.
\bibitem{ref15} NVIDIA, ``Isaac Sim: Robotics Simulation and Synthetic Data,'' 2024.
\bibitem{ref16} M. Jacinto et al., ``Pegasus Simulator,'' \emph{ICUAS}, 2024.
\bibitem{ref17} C. Forster et al., ``SVO: Semi-Direct Monocular Visual Odometry,'' \emph{ICRA}, 2014.
\bibitem{ref18} Y. Yang et al., ``Bio-Inspired Swarm Communication Under Jamming,'' \emph{Drones}, vol.~8, no. 7, 2024.
\end{thebibliography}

\end{document}
